\chapter{Salicylates Level}
\section*{What Is This Test?}
The salicylates level measures the serum concentration of salicylates (primarily acetylsalicylic acid/aspirin, which is rapidly hydrolyzed to salicylic acid after absorption). Salicylate is used therapeutically as an analgesic, antipyretic, anti-inflammatory, and antiplatelet agent. At low doses (81--325 mg/day), aspirin irreversibly inhibits cyclooxygenase-1 (COX-1) in platelets, blocking thromboxane A2 production and platelet aggregation for the prevention of arterial thrombosis (myocardial infarction, stroke). At high doses (3--6 g/day), salicylates are used for anti-inflammatory effects in rheumatic diseases (historical use, largely replaced by other NSAIDs). Salicylate poisoning (acute or chronic) is a potentially life-threatening condition. The classic pathophysiologic triad is: (1) Direct stimulation of the medullary respiratory center causing hyperventilation and primary respiratory alkalosis (earliest sign). (2) Uncoupling of mitochondrial oxidative phosphorylation, producing increased oxygen consumption, increased heat production (hyperthermia), and anaerobic metabolism with lactic acidosis. (3) Metabolic acidosis from accumulation of salicylic acid itself, lactic acid, and ketoacids. The net acid-base disturbance is mixed: primary respiratory alkalosis + increased anion gap metabolic acidosis. CNS toxicity: tinnitus (the earliest and most characteristic symptom of salicylism), hearing loss (reversible, cochlear), and at therapeutic levels, and with toxicity: nausea, vomiting, hyperventilation, confusion, agitation, seizures, coma, cerebral edema, and pulmonary edema (non-cardiogenic pulmonary edema is a cardinal feature of severe salicylate poisoning) \cite{temple1981acute}. Serum salicylate levels are interpreted in conjunction with the clinical presentation using the Done nomogram (for single acute ingestion only --- the Done nomogram should NOT be used for chronic salicylate poisoning, enteric-coated aspirin ingestions, or ingestions of unknown timing) \cite{done1960salicylate}.
\section*{Alternative Names}
Salicylate Level; Aspirin Level; Salicylic Acid Level; Acetylsalicylic Acid Level; ASA Level
\section*{Principle}
Salicylates are measured by a colorimetric method (Trinder reaction) using ferric nitrate reagent, which reacts with salicylic acid to form a violet-colored ferric salicylate complex measured spectrophotometrically at 540 nm. The method is specific for salicylate but can be interfered with by ketones (acetoacetate in diabetic ketoacidosis produces a false-positive color change) and diflunisal (a salicylate-class NSAID). Enzymatic methods (salicylate hydroxylase) are also available on some automated analyzers.
\section*{History}
Salicylic acid, derived from willow bark (Salix), has been used as an antipyretic and analgesic since antiquity; preparations of willow were recommended by Hippocrates in the 5th century BCE. Acetylsalicylic acid was first synthesized by Felix Hoffmann at Bayer in 1897 and was marketed as Aspirin in 1899. John Vane's demonstration that aspirin inhibits prostaglandin synthesis (1971) explained its mechanism of action and earned the 1982 Nobel Prize. The antiplatelet role of low-dose aspirin for prevention of myocardial infarction was established in the 1970s--1980s. Alan Done published the nomogram relating serum salicylate levels to severity of poisoning in 1960, and it guided poisoning management for decades \cite{done1960salicylate}.
\section*{Why This Test Is Ordered}
Diagnosis and management of salicylate poisoning (acute or chronic). In therapeutic use: TDM is NOT routinely indicated for low-dose aspirin antiplatelet therapy. In high-dose aspirin for anti-inflammatory effect (historical), TDM targets 150--300 mcg/mL (mild toxicity), 300--450 mcg/mL (moderate toxicity), >450 mcg/mL (severe toxicity). The Done nomogram correlates serum salicylate level with time since ingestion (for acute single overdose) and predicts severity.
\section*{Primary vs Secondary Test}
This is a secondary, therapeutic-monitoring test rather than a diagnostic test. It is used to confirm and quantitate salicylate exposure in a patient with suspected poisoning and, historically, to guide high-dose anti-inflammatory dosing.
\section*{How This Test Is Performed}
Blood is collected by standard venipuncture into a plain serum tube (red-top) or a serum separator tube (SST, gold-top). For acute poisoning, the level is drawn at least 4--6 hours after ingestion so that absorption is substantially complete, or immediately in the severely ill patient. The serum salicylate concentration is measured by a colorimetric method (the Trinder reaction with ferric nitrate) or by an enzymatic salicylate hydroxylase assay on automated analyzers. Results are reported in mg/dL and are typically available within 1--2 hours for urgent poisoning management.
\section*{What Preparation Is Needed}
No fasting or special preparation is required. The key variable is timing: for the Done nomogram, the level should be drawn at least 4--6 hours after a single acute ingestion, and earlier levels may be falsely low while absorption is ongoing. In chronic poisoning, timing is less useful and the level is drawn whenever the diagnosis is considered, with interpretation based on the clinical picture.
\section*{Contraindications and Precautions}
\begin{itemize}
  \item The Done nomogram applies only to single acute ingestions of immediate-release aspirin; it must not be used for chronic poisoning, enteric-coated preparations, or ingestions of unknown timing
  \item Chronic salicylate poisoning can be fatal at serum levels only mildly above the therapeutic range; interpret the level with the clinical picture, not the nomogram alone \cite{chyka2007salicylate}
  \item Ketones (acetoacetate) and diflunisal interfere with the colorimetric Trinder assay and can falsely elevate results
  \item Acidemia shifts salicylate into tissues and worsens toxicity; correct acidosis even when the measured level is only moderately elevated
  \item Tinnitus and hyperventilation can occur at therapeutic anti-inflammatory levels; treatment is based on the severity of symptoms and the acid-base status
  \item Check for co-ingested acetaminophen in suspected self-poisoning
\end{itemize}
\section*{Risks}
Minimal risk. Slight pain or bruising at the venipuncture site. Rarely: hematoma, infection, or vasovagal syncope.
\section*{What Happens After the Test}
The salicylate level is interpreted against the therapeutic range together with the clinical and acid-base status. Serial levels are often needed in acute poisoning because absorption and metabolism are dynamic and a single early level may understate the peak. Moderate to severe toxicity is treated with IV sodium bicarbonate to alkalinize the urine, potassium repletion, and supportive care; hemodialysis is reserved for severe poisoning, renal failure, or refractory acidemia \cite{omalley2007emergency,proudfoot2004position}. Levels are repeated every 2--4 hours until they are falling consistently and the patient is clinically stable.
\section*{Understanding Results}
Therapeutic: 10--30 mg/dL for analgesic/antipyretic; 150--300 mg/dL for anti-inflammatory. Mild salicylism: 30--50 mg/dL, tinnitus, hearing loss. Moderate toxicity: >50 mg/dL, nausea, vomiting, hyperventilation, respiratory alkalosis. Severe toxicity: >80--100 mg/dL at any time (or >40 mg/dL at 6+ hours): metabolic acidosis, altered mental status, seizures, cerebral edema, non-cardiogenic pulmonary edema, hyperthermia, death. Chronic salicylate poisoning (elderly, renal impairment, therapeutic dosing over days-weeks) is insidious: confusion, delirium, hyperventilation, and metabolic acidosis. Serum level correlates poorly with severity in chronic poisoning; clinical assessment is more important.
\section*{Normal Results}
Therapeutic anti-inflammatory: 150--300 mg/dL (1,500--3,000 mcg/mL). Toxic: >300 mg/dL. The Done nomogram: 40 mg/dL at 6 hours, 30 mg/dL at 12 hours, 20 mg/dL at 24 hours indicate mild toxicity and rising levels above these lines indicate potentially serious poisoning.
\section*{Follow-up Tests}
Arterial or venous blood gas (pH, PCO2, HCO3, mixed acid-base disturbance assessment), serum electrolytes with anion gap, serum creatinine/BUN, serum glucose (CNS hypoglycemia despite normal serum glucose), PT/INR (salicylate interferes with vitamin K-dependent clotting factors), and serum acetaminophen (co-ingestion in suicide attempt).
\section*{References}
\bibliographystyle{unsrtnat}
\bibliography{references}

\clearpage
\section*{Regulatory Status and Authorization Date}
This chapter describes a test category, not a specific commercial product. FDA, CDSCO, EU IVDR, and MHRA approval or registration dates therefore cannot be assigned without the assay manufacturer, product name, instrument, and intended use. For laboratory-developed tests, an individual agency approval date may not exist. See the Preface for the interpretation of regulatory status.

\clearpage
